\section{Open Drain, and the Line That Has Two Masters}
\label{sec:gpio-open-drain}

An I2C bus has two wires and, on each of them, two possible masters, each of which may want to pull
the line low while the other wants it high. The electrical answer is the reason the mode list has a
\codeinline{WIRED\_AND} in it.

\subsection{What WIRED\_AND Does}

\begin{minted}{text}
mode = WIRED_AND   (8)
  DOUT bit = 0  ->  the pad pulls the line to ground
  DOUT bit = 1  ->  the pad releases the line (high impedance)
\end{minted}

The line is held high by an external resistor. So ``1'' is not something the pin produces — it is
something the pin \emph{stops preventing}. Two devices can then share the wire: any one of them can
pull it low, and the only way it goes high is if nobody does. That is what makes arbitration
possible, and it is why a push-pull pin must never be connected to a shared bus: driving it high
against another device driving low is a short circuit through two transistors.

\subsection{The Measurement That Proved the Pull-Ups}

The IMU could not be found for two sessions. Before blaming the pin map, this project measured
whether the bus lines had pull-ups at all:

\begin{minted}{text}
# every pin on PORTB set to input, then read DIN
PORTB_MODEL = 0x11111111     # all inputs
PORTB_DIN   = 0x0000000C     # bits 2 and 3 high
\end{minted}

Bits 2 and 3 are \codeinline{PB2} and \codeinline{PB3}: the two pins of the IMU's bus, reading high
with nothing driving them. That is an external pull-up, present and working — which retired the
``the board has no pull-ups'' hypothesis in one instruction and moved the investigation to the pin
\emph{order}.

\subsection{Bit-Banging an Open-Drain Bus}

With \codeinline{WIRED\_AND} the whole I2C protocol reduces to two writes and one read:

\begin{minted}{rust}
fn sda(&self, high: bool) { write_dout(self.gpio, self.bus.port, self.bus.sda_bit, high) }
fn scl(&self, high: bool) { write_dout(self.gpio, self.bus.port, self.bus.scl_bit, high) }
fn read_sda(&self) -> bool { read_din(self.gpio, self.bus.port) & (1 << self.bus.sda_bit) != 0 }
\end{minted}

\textbf{There is no direction switch.} On an open-drain bus a pin used as an input with the output
latch high is the same electrical state as a released output — which is why a bit-banged I2C master
can be four functions and no state machine, and why the same pin can be written and read in the same
transfer.

\subsection{Half a Bit Period, and Why It Is a Number in Cycles}

Timing on a bit-banged bus is a delay per half period, and the delay is not a duration — it is a
count of loop iterations whose real length depends on the compiler:

\begin{minted}{rust}
// ~5 us at 19 MHz, measured: 60 iterations of this loop on this compiler.
const HALF_BIT_CYCLES: u32 = 60;
\end{minted}

That constant is a measurement, not a datasheet value, and it is the reason chapter~14 exists: the
same class of mistake — \codeinline{asm::delay(n)} assumed to be \codeinline{n} cycles when it is
about three per iteration — turned a 2.7 kHz tone into something else entirely.
